% !TeX program = lualatex
\documentclass[10pt]{ltjsarticle}
\usepackage[margin=1in]{geometry}

% Core packages
\usepackage{amsmath,amssymb}
\usepackage{tikz-cd}
\usepackage{multicol}
\usepackage{enumitem}
\usepackage{hyperref}

% Paragraphs
\setlength{\parindent}{0pt}
\setlength{\parskip}{1\baselineskip}

\title{時間付きブール代数とその一階述語論理}
\author{千葉将臣}
\date{2026-07-13}

\begin{document}
\maketitle

\section{時間付きブール代数の基本アイデア}

\subsection{時間の一方向性}

時間は過去から未来へ一方向にのみ流れるものとする。各時間点 $t$ において、命題 $A$ はブール値
\[
  A(t) \in \{0,1\}
\]
を持つ。

\subsection{否定の時間的解釈}

本稿では、否定記号 $\neg$ を通常の静的な反転ではなく、未来への到達可能性を含む演算として解釈する。すなわち、現在時刻を $\mathsf{now}$ とすると、
\begin{align*}
  \neg A &:= \text{将来いつか $A$ が偽になる}, \\
  \neg\neg A &:= \text{将来いつか $A$ が真になる}, \\
  \neg\neg\neg A &:= \text{将来いつか $A$ が偽になる}
\end{align*}
と読む。この解釈では、否定を反復することにより、真偽値そのものではなく、命題の時間的変化可能性が表現される。

\subsection{\texorpdfstring{矛盾位相 $A \land \neg A$}{矛盾位相 A and not A}}

古典論理において $A \land \neg A$ は矛盾であり恒偽である。しかし、時間付きブール代数では
\[
  A \land \neg A
  \quad\text{を}\quad
  \text{今 $A$ が真であり、かつ将来いつか $A$ が偽になる}
\]
と解釈する。したがって、時間的に変化する命題に対しては、これは単なる矛盾ではなく、現在と未来の差異を表す位相である。

例えば「今は晴れているが、将来いつか雨が降る」という命題は、静的には矛盾に見える形式を持ちながら、時間発展を含めれば自然に理解できる。

\subsection{\texorpdfstring{独自の公理 $\neg\neg\neg A \to (\neg A \land \neg\neg A)$}{独自の公理 not not not A implies not A and not not A}}

次の公理を考える。
\[
  \neg\neg\neg A \to (\neg A \land \neg\neg A).
\]
ここで $\neg\neg\neg A$ は「将来いつか $A$ が偽になる」ことを表し、$\neg A \land \neg\neg A$ は「将来いつか $A$ が偽になり、かつ将来いつか $A$ が真になる」ことを表す。

したがって、この公理は次の意味を持つ。
\[
  \text{$A$ が将来偽になるならば、$A$ は真と偽の間を往復する。}
\]
すなわち、この体系では、否定の反復は単なる古典的な二重否定消去ではなく、命題の振動的な時間発展を記述する。

\section{FOLでの構成}

\subsection{時間を変数として導入する}

時間 $t$ を一次の変数として導入する。時間領域は自然数、整数、実数、あるいは任意の線形順序集合としてよい。命題 $A$ は時間に依存する述語 $A(t)$ として表現される。

\subsection{否定のFOL表現}

現在時刻を $\mathsf{now}$ と書く。このとき、時間付き否定は一次論理において次のように表現できる。
\begin{align*}
  \neg A &\equiv \exists t\,(t > \mathsf{now} \land \neg A(t)), \\
  \neg\neg A &\equiv \exists t\,(t > \mathsf{now} \land A(t)), \\
  \neg\neg\neg A &\equiv \exists t\,(t > \mathsf{now} \land \neg A(t)).
\end{align*}
この定義では、否定は未来量化と古典的否定の組として表される。

\subsection{矛盾位相のFOL表現}

時間付きブール代数における $A \land \neg A$ は、FOLでは
\[
  A(\mathsf{now}) \land \exists t\,(t > \mathsf{now} \land \neg A(t))
\]
と表される。これは「今 $A$ が真であり、将来いつか $A$ が偽になる」という意味であり、時間変化を許すモデルでは充足可能である。

\subsection{独自公理のFOL表現}

公理
\[
  \neg\neg\neg A \to (\neg A \land \neg\neg A)
\]
は、FOLでは次のように表現できる。
\[
  \bigl(\exists t_1\,(t_1 > \mathsf{now} \land \neg A(t_1))\bigr)
  \to
  \bigl(
    \exists t_2\,(t_2 > \mathsf{now} \land \neg A(t_2))
    \land
    \exists t_3\,(t_3 > \mathsf{now} \land A(t_3))
  \bigr).
\]
これは「将来いつか $A$ が偽になるならば、将来いつか $A$ が偽になり、かつ将来いつか $A$ が真になる」という振動性の条件である。

\section{プログラミング言語との対応}

\subsection{プログラミングを代数にしたいという背景}

プログラミングを代数化しようとする試みは、計算機科学の中で繰り返し現れてきた。例えば、プログラム代数ではプログラムを項として扱い、合成・選択・反復などの演算を代数法則によって記述する。プロセス代数では、並行プロセスを代数的対象として扱い、通信・同期・相互作用を演算として表現する。代表例として CCS、CSP、$\pi$ 計算などがある。また、ホーア論理や動的論理は、プログラムの実行前後に成り立つ性質を論理式として記述し、検証可能にする枠組みである。

これらに共通しているのは、第一にプログラムの構造を演算として捉えること、第二にプログラムの性質を論理的に記述することである。すなわち、プログラミングを代数化するとは、単にコードを記号列として扱うことではなく、プログラムの合成、分岐、反復、通信、状態変化を、計算可能な代数的構造として整理することである。

時間付きブール代数は、この流れに対して、状態の真理値とその時間発展に焦点を当てた別の入口を与える。プログラムの実行を状態の列、あるいは分岐を許す状態の木として見れば、各状態で命題 $A$ は真または偽の値を持つ。この状況は、各時間点 $t$ において $A(t) \in \{0,1\}$ を考える時間付きブール代数と自然に対応する。

\subsection{時相論理・動的論理との関係}

時相論理では、$F A$ は「将来いつか $A$ が真になる」、$G A$ は「今後ずっと $A$ が真である」ことを表す。動的論理では、$[\alpha]P$ は「プログラム $\alpha$ の実行後、必ず $P$ が成り立つ」ことを表す。

時間付きブール代数の $\neg A$ を
\[
  F(\neg A)
\]
と解釈すれば、この体系は時相論理の枠組みに回収できる。すなわち、時間付き否定は、古典的否定と未来様相 $F$ の合成として理解できる。

\subsection{否定だけで時間性が表現できる、という意味}

通常の時相論理では、時間性を表すために $F, G, U$ などの専用演算子を導入する。これに対して、時間付きブール代数では、否定 $\neg$ だけで時間的な様相を表現する。
\begin{align*}
  \neg A &:= \text{将来いつか $A$ が偽になる}, \\
  \neg\neg A &:= \text{将来いつか $A$ が真になる}.
\end{align*}
これは、否定という基本的な論理演算に、時間性という豊かな構造が埋め込まれていることを意味する。

\subsection{状態遷移としての否定}

プログラムの実行を状態遷移系として見ると、時間付きブール代数の否定 $\neg$ は単なる真偽値の反転ではなく、状態変化や実行ステップを表す演算として読める。すなわち、
\[
  \neg A := \text{将来いつか $A$ が偽になる}
\]
という定義は、現在の状態から到達可能な未来状態において $A$ が破れることを意味する。これは、プログラムの実行によって性質 $A$ が保存されるか、あるいは将来失われるかを記述している。

同様に、
\[
  \neg\neg A := \text{将来いつか $A$ が真になる}
\]
は、現在 $A$ が成り立っていない場合でも、実行の進行によって $A$ が再び成立する可能性を表す。この意味で、否定の反復は、真偽値の静的な反転ではなく、実行の時間的深さを表現する。

したがって、時間付きブール代数では、否定という最も基本的な論理演算の内部に、プログラムの実行ステップや状態変化が埋め込まれていると解釈できる。

\subsection{矛盾位相と変化の過程}

時間付きブール代数における
\[
  A \land \neg A
\]
は、古典論理における矛盾ではなく、
\[
  \text{今 $A$ が真であり、将来いつか $A$ が偽になる}
\]
という状態変化の過程を表す。プログラム意味論の観点からは、これは、ある性質 $A$ が現在の状態では成り立っているが、実行を進めると破れる可能性があることを意味する。

例えば、$A$ を「変数 $x$ は正である」という命題とする。このとき $A \land \neg A$ は、「現在は $x>0$ であるが、将来の実行状態では $x \leq 0$ になりうる」という意味になる。これは矛盾ではなく、プログラムの状態更新によって性質が変化することの表現である。

このように、矛盾位相は、静的論理では排除される形式を、動的な変化の記述として再解釈する役割を持つ。

\subsection{振動公理とループ的挙動}

独自公理
\[
  \neg\neg\neg A \to (\neg A \land \neg\neg A)
\]
は、$A$ が将来偽になるならば、将来偽になる局面と真になる局面の両方が存在することを要求する。プログラムの振る舞いとして読むならば、これは状態が $A$ の成立と不成立の間を往復することを意味する。

したがって、この公理はループや反復的な挙動を記述するための抽象的な原理として解釈できる。例えば、あるフラグが実行中にオンとオフを繰り返す、ある資源が獲得と解放を繰り返す、ある条件が周期的に成立と不成立を行き来する、といった状況を表現できる。

もちろん、すべてのプログラムがこの振動性を持つわけではない。したがって、この公理は一般のプログラム意味論に常に課す法則というより、反復・周期・再帰的挙動を持つ計算を特徴づける追加公理として位置づけるのが自然である。

\subsection{既存のプログラム代数との比較}

プロセス代数は、プロセスそのものを項として扱い、並行合成、選択、通信、同期などを演算として表現する。これに対して、時間付きブール代数は、プロセスの外形よりも、各状態で成り立つ命題の真理値とその変化に焦点を当てる。したがって、プロセス代数が「何が実行されるか」を代数化する枠組みであるのに対し、時間付きブール代数は「実行によって何が真または偽になるか」を代数化する枠組みである。

動的論理との関係では、プログラム $\alpha$ の実行後の性質を $[\alpha]P$ や $\langle\alpha\rangle P$ によって表すのに対し、時間付きブール代数では未来への到達可能性を否定 $\neg$ の中に埋め込む。このため、時間付きブール代数は、動的論理の様相演算子を、否定演算の時間的解釈として圧縮した体系と見ることができる。

この意味で、時間付きブール代数は、プログラム代数やプロセス代数を置き換えるものではなく、それらを状態命題の時間的変化という側面から補完する。両者を組み合わせれば、プロセスの構造と状態の真理値変化の双方を代数的に扱うことができる。

\subsection{モデル検査への応用}

プログラムやシステムの振る舞いを状態遷移系としてモデル化する。各状態で命題が真または偽のブール値を持つとすれば、時間付きブール代数によってシステムの性質を記述できる。

例えば「ずっと安全である」「いつか終了する」といった性質は、時相論理やFOLを通じて形式化され、モデル検査によって自動検証の対象となる。

以上より、時間付きブール代数によって、プログラムの状態遷移を時間付きブール値として表現できる。また、否定 $\neg$ を実行ステップや状態変化として読むことで、矛盾位相 $A \land \neg A$ は変化の過程を、振動公理 $\neg\neg\neg A \to (\neg A \land \neg\neg A)$ はループ的・反復的な振る舞いを記述する。FOLで構成すれば、時相論理や動的論理と接続できるため、プログラミング言語の意味論や検証への応用が期待できる。

\section{圏論的な視点}

\subsection{モナドとしての時間}

通常のブール代数の圏では、対象はブール代数、射はブール準同型である。時間付きブール代数では、各ブール値 $b$ を時間発展したブール値の列 $T(b)$ に写すモナド $T$ を考える。

このとき、単位
\[
  \eta : b \to T(b)
\]
は、現在のブール値を「ずっと同じ値を取る列」として埋め込む操作である。また、乗法
\[
  \mu : T^2(b) \to T(b)
\]
は、「将来の将来」を一つの将来へと平坦化する操作である。

\subsection{FOLとの対応}

FOLで時間量化 $\exists t, \forall t$ を行うことは、モナド $T$ を時間量化の構造として解釈することに対応する。この観点から、時間付きブール代数は、プログラミング言語における IO、State、Maybe などのモナドと類似した意味論的構造を持つ。

\section{哲学的な解釈}

\subsection{否定と時間の深い関係}

否定は「$A$ ではない」という差異を表す。一方で、時間は「今と未来」あるいは「過去と現在」という差異を生み出す。

時間付きブール代数では、否定が時間的な差異を駆動する。
\begin{align*}
  \neg A &:= \text{将来いつか $A$ が偽になる}, \\
  \neg\neg A &:= \text{将来いつか $A$ が真になる}.
\end{align*}
したがって、否定は単なる補集合操作ではなく、未来への変化可能性を含んだ演算として理解される。

\subsection{論理主義的な時間観}

この体系は、「世界の変化は、否定という論理的操作から生じる」という解釈を与える。時間付きブール代数は、この考え方を数学的に定式化する試みと見なせる。

\section{まとめ}


\subsection{実現できること}

以上の議論をまとめると、時間付きブール代数によって実現できることは次のように整理できる。

\begin{enumerate}[label=\arabic*.]
  \item \textbf{否定だけによる時間性の表現}\\
  通常の時相論理では $F, G, U$ などの専用の時間演算子を導入するが、時間付きブール代数では否定 $\neg$ を「将来いつか反対の値になる」という操作として解釈することで、否定だけから時間的な変化可能性を表現できる。

  \item \textbf{変化の過程としての矛盾位相}\\
  $A \land \neg A$ は古典論理では矛盾であるが、この体系では「今 $A$ が真であり、将来いつか $A$ が偽になる」という変化の局面を表す。したがって、静的には矛盾に見える形式を、時間発展を含む状態記述として扱える。

  \item \textbf{振動・往復する振る舞いの記述}\\
  公理 $\neg\neg\neg A \to (\neg A \land \neg\neg A)$ は、$A$ が将来偽になるならば、将来偽になる局面と真になる局面の両方が存在することを要求する。これは、ループ、周期的変化、フラグのオン・オフ、資源の獲得と解放のような往復的挙動を抽象的に記述するための原理になる。

  \item \textbf{FOLによる形式化}\\
  時間を一次の変数 $t$ として導入し、命題を時間依存述語 $A(t)$ と見れば、時間付き否定は $\exists t\,(t > \mathsf{now} \land \neg A(t))$ のような一次論理式として表現できる。これにより、独自の記法をFOLの枠組みに翻訳し、通常の論理的推論やモデル理論と接続できる。

  \item \textbf{プログラムの状態遷移の代数的記述}\\
  プログラムの実行を状態の列または木と見れば、各状態で命題 $A$ が真偽値を持つ。時間付きブール代数は、この真偽値の時間発展を代数的に扱うため、実行ステップ、状態更新、性質の保存や破壊を簡潔に記述できる。

  \item \textbf{検証・モデル検査への応用}\\
  $\neg A$ を $F(\neg A)$ と読むことで、時間付きブール代数は時相論理と接続する。したがって、「いつか終了する」「将来エラーになる可能性がある」「安全性が破れる可能性がある」といった性質を、状態遷移系上で検査するための記述言語として利用できる。

  \item \textbf{モナドとしての時間・計算効果の表現}\\
  ブール値を時間発展する列へ写す操作をモナド $T$ と見れば、現在値の埋め込み $\eta$ と「将来の将来」の平坦化 $\mu$ によって、時間や計算効果を圏論的に整理できる。この意味で、時間付きブール代数は IO、State、Maybe などのプログラミング言語のモナドと類似した構造を持つ。

  \item \textbf{否定と時間の関係の哲学的定式化}\\
  否定は「$A$ ではない」という差異を表し、時間は「今と未来」の差異を生み出す。時間付きブール代数では、否定が時間的差異を生成する操作として働くため、「変化は否定から生じる」という論理主義的な時間観を数学的に表現できる。
\end{enumerate}

時間付きブール代数は、否定 $\neg$ だけで時間性を表現する独自の論理体系である。FOLで構成すると、時相論理や動的論理の枠組みを通じて、プログラミング言語の意味論や検証と密接に結びつく。

圏論的には、時間や計算をモナドとして表現する視点を与える。哲学的には、否定と時間の深い関係を示唆する、論理主義的な時間観として興味深い。

\end{document}
